% Propietario conservado: /Users/ruben/Documents/New project/output/REVISION_ARGUMENTAL_APERTURAS_CIERRES_20260915/VI/delivery/MANUSCRIPT_VI_ES_EN_COMPLETE/source_es/sections/02_particula_persistencia.tex
\section{La particule comme classe de sections persistantes}
\label{vi:vi:sec:particula}

La construction matérielle commence dans l’état produit par APP, TRIT et
TPK. APP fournit les positions et les deux opérations arithmétiques ;
TRIT distingue le régime et l’orientation de leurs transitions ; TPK
transporte cette information avec les quotients, les incidences de
frontière et la mémoire. Lors du prolongement de l’état, une évaluation finie
reste liée aux évaluations qui la précèdent. La notion de
particule s’appuie sur cette compatibilité et sur sa représentation matérielle.
Sa masse sera une observable de la section ainsi construite.

Cette séquence détermine le sens des trois objets qui seront
utilisés : l’état généalogique conserve l’extension complète ; la route
enregistre son histoire observable ; la représentation matérielle fait agir
cette histoire sur un espace d’états et permet d’évaluer des observables.
La distinction est essentielle même lorsque les trois objets partagent
une même étiquette électronique ou une même coordonnée de l’atlas.

\subsection{Extension, histoire observable et section compatible}
\label{vi:vi:subsec:extension-persistente}

Soient $\mathcal S_m$ les états admissibles survivants jusqu’à la profondeur
$m$, obtenus par le prolongement du TPK, et soient
$\rho_{n,m}:\mathcal S_n\to\mathcal S_m$, pour $n\geq m$, leurs
applications de restriction. La compatibilité du raffinement signifie
\begin{equation}
 \rho_{m,m}=I,\qquad
 \rho_{n,m}\rho_{p,n}=\rho_{p,m}\quad(p\geq n\geq m).
 \label{vi:vi:eq:sistema-estados}
\end{equation}
Une extension survivante est un élément
\begin{equation}
 x=(x_m)_{m\geq0}\in\Omega^{\mathrm{surv}}_{\mathrm{HMT}}
 :=\varprojlim_m\mathcal S_m,
 \qquad \rho_{n,m}(x_n)=x_m.
 \label{vi:vi:eq:extension-superviviente}
\end{equation}
L’existence et l’admissibilité de ces états proviennent de la
construction du noyau formel. L’expression de limite projective décrit
les compatibilités que conserve un état déjà engendré.

Soit $\mathcal R_m$ l’espace des histoires observables à la profondeur $m$.
Le lecteur de route $\operatorname{sh}_m:\mathcal S_m\to\mathcal R_m$
et la restriction $\operatorname{tr}_{n,m}$ satisfont
\begin{equation}
 \operatorname{tr}_{n,m}\operatorname{sh}_n
 =\operatorname{sh}_m\rho_{n,m}.
 \label{vi:vi:eq:naturalidad-ruta}
\end{equation}
Ainsi, le lecteur transmet les restrictions de l’état à ses histoires.
La route de $x$ est la famille
$\gamma_x=(\gamma_x^{[m]})_m$, où
$\gamma_x^{[m]}=\operatorname{sh}_m(x_m)$.

\begin{proposition}[Compatibilité de l’histoire publiée]
\label{vi:vi:prop:historia-compatible}
Pour toute extension \eqref{vi:vi:eq:extension-superviviente}, ses histoires
finies satisfont
$\operatorname{tr}_{n,m}(\gamma_x^{[n]})=\gamma_x^{[m]}$.
\end{proposition}
\begin{proof}
En appliquant \eqref{vi:vi:eq:naturalidad-ruta} à l’état $x_n$, on obtient
\[
 \operatorname{tr}_{n,m}(\gamma_x^{[n]})
 =\operatorname{sh}_m(\rho_{n,m}x_n)
 =\operatorname{sh}_m(x_m)=\gamma_x^{[m]}.
\]
La seconde égalité utilise précisément la condition de limite projective.
\end{proof}

La proposition établit la descente de l’extension à l’histoire.
Son contenu n’inclut pas l’injectivité du lecteur : deux extensions
peuvent publier une même histoire finie et conserver des mémoires distinctes.
Ainsi, lorsqu’une route est réutilisée, son état de provenance est conservé,
au lieu de le reconstruire à partir d’un chiffre ou d’un mot visible.

Sur chaque histoire $\gamma_x^{[m]}$, on considère le fibré discret
$\mathcal B_{\gamma_x^{[m]}}$ des états matériels de la représentation
déclarée. Soient
$\mathscr E_m(x)=\Gamma(\mathcal B_{\gamma_x^{[m]}})$ ses espaces
de sections, avec des restrictions
$p_{n,m}:\mathscr E_n(x)\to\mathscr E_m(x)$ compatibles avec celles
de l’état. Une section persistante est une famille
\begin{equation}
 s=(s_m)_{m\geq m_0},\qquad s_m\in\mathscr E_m(x),\qquad
 p_{n,m}s_n=s_m\quad(n\geq m\geq m_0).
 \label{vi:vi:eq:seccion-persistente}
\end{equation}
La profondeur initiale $m_0$ permet de décrire un état depuis toute
coupure à partir de laquelle sa représentation est définie. L’information
antérieure qui intervient dans ses observables demeure dans le registre
transporté par l’extension.

\subsection{Identité sous raffinement et équivalence des représentations}
\label{vi:vi:subsec:identidad-material}

Fixons les groupes de symétries internes $G_m$ qui agissent sur les fibres
matérielles et préservent les observables déclarées. Deux sections
persistantes $s$ et $s'$ représentent la même identité matérielle lorsque,
à partir d’un certain niveau $m_1$, il existe une famille $g_m\in G_m$ telle que
\begin{equation}
 s'_m=g_ms_m,\qquad
 p_{n,m}g_n=g_mp_{n,m},
 \label{vi:vi:eq:equivalencia-persistente}
\end{equation}
et que les caractères observables du registre généalogique coïncident. Les
symétries admissibles sont fixées avant de former le quotient ; leur action
transporte le feuillet, l’orientation, la mémoire et la représentation avec
les lois correspondantes.

\begin{proposition}[Équivalence compatible des sections]
\label{vi:vi:prop:equivalencia-persistente}
La relation \eqref{vi:vi:eq:equivalencia-persistente}, sur des familles qui
portent les mêmes types de représentation et d’observables, est une relation
d’équivalence. Ses classes ne dépendent pas du retrait d’un nombre
fini de niveaux initiaux, pourvu que leur information
généalogique transportée soit conservée.
\end{proposition}
\begin{proof}
L’identité de chaque $G_m$ donne la réflexivité. L’inverse d’une famille
compatible est encore compatible : de
$p_{n,m}g_n=g_mp_{n,m}$, on déduit
$p_{n,m}g_n^{-1}=g_m^{-1}p_{n,m}$. Cela donne la symétrie.
Si $s'=gs$ et $s''=hs'$, à partir du plus grand des deux niveaux initiaux,
on a $s''_m=h_mg_ms_m$ et
$p_{n,m}h_ng_n=h_mg_mp_{n,m}$. La composition préserve également les
observables et prouve la transitivité. Toutes ces relations sont formulées
sur une queue commune ; changer le point initial dans cette queue
conserve la classe. La dernière assertion concerne le changement de
présentation, et non l’effacement de la mémoire de l’état.
\end{proof}

\begin{definition}[Particule matérielle persistante]
\label{vi:vi:def:particula}
Une particule HMT--MD est une classe $[s]_{\mathrm{pers}}$ de sections
\eqref{vi:vi:eq:seccion-persistente} sous
\eqref{vi:vi:eq:equivalencia-persistente}, accompagnée de son extension et de son
registre de provenance, qui admet une représentation matérielle
irréductible ou un bloc minimal invariant et une sortie spectrale stable
sous le couplage spécifié.
\end{definition}

La définition réunit des conditions de différents niveaux. Le prolongement
et le registre appartiennent à l’état HMT. La représentation et le
couplage appartiennent à sa réalisation matérielle. La stabilité
spectrale se rapporte à l’opérateur et au domaine utilisés par ce
couplage. Cette organisation permet d’étudier conjointement identité,
charge, spin et masse sans transformer une observable particulière en définition
de toutes les autres.

Le retour nonadique illustre la distinction entre identité et état
instantané. Pour le lecteur de phase $g$ et le registre de mémoire
$\mathcal B$, peuvent coexister
\begin{equation}
 g(k+9)=g(k),\qquad
 p_{k+9,k}(s_{k+9})=s_k,\qquad
 \mathcal B_{k+9}\ne\mathcal B_k.
 \label{vi:vi:eq:persistencia-fase-memoria}
\end{equation}
La première égalité est périodique ; la seconde exprime la persistance ; la
troisième enregistre l’avancée. Le retour de phase appartient à l’histoire
d’une même identité matérielle, tandis que l’état conserve le fait qu’un
tour supplémentaire a été effectué.

\subsection{L’espèce comme fibre : portée de l’atlas fini}
\label{vi:vi:subsec:especie-fibra}

L’atlas construit par les opérations APP--TRIT--TPK a pour support
$\mathcal C_{81}=\{1,\ldots,9\}^2$. Ses quatre orientations par
cellule forment le domaine de $324$ routes orientées. L’application de
registre $\mathcal R$ distingue $56$ familles de routes, tandis que le
classificateur interne
\begin{equation}
 q:\mathcal C_{81}\longrightarrow\Sigma_{13}
 \label{vi:vi:eq:clasificador-especies}
\end{equation}
réunit $13$ types de famille et de niveau. Ces recensements et leurs règles sont
développés dans le corps consacré à l’atlas ; on précise ici quel objet
ils transmettent à la définition de particule.

Pour $y\in\Sigma_{13}$, la fibre $q^{-1}(y)$ conserve toutes ses
cellules et tous ses registres. L’ordre interne de chaque fibre de
$\mathcal R$ étant fixé, soit $k(c)$ le rang de la cellule au sein de celle-ci. Sa
présentation finie est
\begin{equation}
 \mathfrak S_y^{\mathrm{atlas}}
 =\{(\mathcal R(c),k(c)):q(c)=y\}.
 \label{vi:vi:eq:multiseccion-atlas}
\end{equation}
Le rang est ajouté pour séparer les cellules qui possèdent le même registre de
famille ; il ne sélectionne pas une masse. La paire $(\mathcal R(c),k(c))$
retrouve la cellule dans la carte ordonnée, car la seconde donnée
individualise sa position dans la fibre de la première.

Si $\pi_{81}:\Omega^{\mathrm{surv}}_{\mathrm{HMT}}\to
\mathcal C_{81}$ est le lecteur cellulaire, le relèvement de l’espèce est
\begin{equation}
 \widetilde{\mathfrak S}_y
 =\{x\in\Omega^{\mathrm{surv}}_{\mathrm{HMT}}:
                  q(\pi_{81}x)=y\}.
 \label{vi:vi:eq:multiseccion-enriquecida}
\end{equation}
Cette préimage contient des extensions, tandis que
\eqref{vi:vi:eq:multiseccion-atlas} contient des présentations finies. La
construction et le domaine du lecteur $\pi_{81}$ font partie du passage
entre les deux. Une particule concrète se réalise comme section persistante
au sein de la multisection enrichie correspondante.

\begin{theorem}[Point fixe central et sélection équivariante]
\label{vi:vi:thm:selector-central}
Soit $\rho_c(r,s)=(10-r,10-s)$ l’inversion cellulaire. Dans les fibres
$\rho_c$-stables du classificateur \eqref{vi:vi:eq:clasificador-especies},
avec une action triviale sur l’étiquette $y$, une sélection ponctuelle
équivariante ne peut prendre de valeurs que dans la cellule $(5,5)$. La fibre
électronique de niveau zéro contient cet unique point fixe ; aucune autre
fibre n’admet une telle sélection.
\end{theorem}
\begin{proof}
Résoudre $\rho_c(r,s)=(r,s)$ donne $r=s=5$. Si $a(y)\in q^{-1}(y)$
est une sélection équivariante et si $y$ reste fixe, alors
$a(y)=a(\rho_c y)=\rho_c a(y)$. Sa valeur doit être le point fixe
$(5,5)$. L’atlas situe cette cellule dans la fibre électronique de niveau
zéro, dont le cardinal vaut un. Pour toute autre fibre, la condition
de sélection produirait un point fixe qui ne lui appartient pas.
\end{proof}

La conséquence est constructive : les autres espèces sont présentées par
orbites ou multisections, et une réalisation qui nécessite une droite
particulière doit déclarer sa représentation et son couplage.
L’électron fournit le cas central de rang un sans transformer toutes
les familles en copies de ce cas. De plus, la cellule $(5,5)$, la signature
centrale brute, la signature régionale $555555$ et l’origine relative d’une
droite d’action conservent leurs domaines propres. L’unicité du point
n’identifie pas ces registres entre eux.

\subsection{Construction de l’observable de masse à partir du registre matériel}
\label{vi:vi:subsec:particula-observable}

La chaîne utilisée par l’observable de masse conserve l’ordre
\begin{equation}
 \text{extension}\longrightarrow\text{route}\longrightarrow
 \text{registre}\longrightarrow\text{signature entière}\longrightarrow
 \text{caractère}\longrightarrow\text{tours}\longrightarrow
 \text{opérateur matériel}.
 \label{vi:vi:eq:cadena-material}
\end{equation}
Chaque flèche retient les données nécessaires à la suivante. La signature est
une lecture du registre ; le caractère évalue cette signature avec les
coefficients internes déjà construits ; les tours incorporent son
raffinement ; l’opérateur agit finalement dans la représentation de la
fibre sélectionnée par le couplage. Sa sortie peut être un scalaire,
un multiplet ou une distribution spectrale, selon cette représentation.
La comparaison métrologique a lieu après cette construction.

Une espèce est décrite par sa charge, sa représentation de spin, sa couleur,
sa saveur, sa génération, sa famille de routes et son opérateur matériel. Ces données
expliquent pourquoi deux préparations de masse coïncidente peuvent
conserver des informations distinctes, et pourquoi une même identité peut
publier des observables différentes sous des couplages distincts.
L’identité réside dans la classe persistante et dans ses transports admissibles ;
la mesure est une lecture de cette identité dans le codomaine déclaré.

La même hiérarchie sépare une extension persistante d’une section
d’horizon fini. Si $\operatorname{Ext}_L(s)\ne\varnothing$ et
$\operatorname{Ext}_{L+1}(s)=\varnothing$, l’horizon $L$ enregistre
une obstruction au prolongement. Une participation intermédiaire de ce
type dans une composition ne crée pas, à elle seule, une nouvelle espèce
asymptotique. De même, une résonance exige l’opérateur de
survie et son pôle, en plus de l’identité de la section qu’elle
transporte. Ainsi est préservée la différence entre persistance,
prolongement fini et évaluation spectrale.
